\chapter{如何寫小說：我親自驗證過的 12 步流程}

這場講座雖然沒有黑板上的推導，卻有一種在結構上極為接近推導的東西。Jenkins 把初學者對「寫小說」那種模糊而擴散的恐懼，逐步轉化成一連串限制、測試與更新規則。這個順序很重要。我們先界定對象，再選定計畫，再決定它要如何被承載、如何被觀看、如何開場，以及如何一路被推向最大壓力點，最後才讓它獲得解決。

\begin{remark}
本章中的顯示公式，都是對講座中程序性邏輯所做的審慎重建。它們不是從畫面上的記號逐字轉錄而來，而只是把建議的結構明白地攤開。
\end{remark}

\section{從恐懼走向方法}

Jenkins 一開始先清理地面。先於流程、先於技藝、也先於動機，他先告訴我們：這裡討論的究竟是什麼對象。在這場講座裡，所謂小說，就是虛構作品。這聽起來像是最基本不過的事，但這正是恰當的第一步，因為它固定了目標，也防止整場談話滑散成對「寫作」這件事過於寬泛的討論。

\begin{definition}
就本講座的意義而言，小說是一種虛構作品。
\end{definition}

\begin{equation}
\text{小說}=\text{虛構}.
\end{equation}

從這裡開始，講座立刻轉向真正的障礙。聽眾心中有一個想法，想寫作，甚至也可能早就被人說過自己有文字天分；但私底下，總有一種恐懼：一整本小說會不會大到寫不完，或者難到根本不可能寫得好。Jenkins 並不以模糊的「靈感」來安撫這種恐懼，而是以一套可重複的方法來作答。

同時，他也插入了兩個值得保留的動機性節點。第一，即便是有經驗的小說家，也依然會擔心一個點子是否撐得起一本到書，以及自己的文字是否夠好。第二，他要求聽眾把整場講座看到最後，因為最後還會有一個關於修訂的 bonus step，以及一份在投稿前逐字檢查的 checklist。因此，講座從一開始就在做它自己所主張的事：先立出問題，再許諾一套結構，然後預先設下一個回報點。

\subsection{Question \& Answer}

\paragraph{Question.} 如果我們害怕自己根本沒有能力寫完一整部小說，該怎麼辦？

\paragraph{Answer.} 我們不等恐懼自行消失，而是用一套具體步驟取代一種模糊的野心。Jenkins 第一個真正的承諾，不是寫作會變得容易，而是它會變得可理解、可重複。

\section{選故事，以及如何熬過中段}

第一個真正的實作問題，是選擇。我們可能手上有很多個故事點子；Jenkins 把這件事當成初始條件，而不是解答本身。真正的問題，是怎樣從中選出那一個：它既有足夠重量撐起一本到書，也有足夠情感力量，讓我們在最初幾頁的新鮮感消退之後，仍然願意繼續工作。

他立刻給出一個實用尺度：一個真正能勝出的點子，必須能支撐一部相當長度的小說，而不是只支撐一個聰明的開頭、一場戲，或一個前提。從結構上看，小說首先會被分成三個區域：

\begin{equation}
\text{開場} \to \text{中段} \to \text{結尾}.
\end{equation}

這一節中最令人難忘的一個說法，就是用來命名那塊最困難區域的詞。

\begin{definition}
所謂「中段馬拉松」，是指小說介於開場與結尾之間那段漫長的核心區域，在那裡，戲劇、張力、衝突、行動，以及所有等待回收的鋪陳，都必須持續被維持。
\end{definition}

在這裡，Jenkins 比一般那種 ``寫你所愛'' 的建議更精確。光是在一開始讓人興奮還不夠；一個點子還必須能熬得住時間。測試一個點子，不能只看它是否新穎，還要看：當開場的能量已經退去，而結尾又還遙不可見時，它是否仍然會把我們拉回去。

接著，講座給出了一個異常有用的反應測試。Jenkins 把作者自己的反應，當成頁面上故事能量的一種局部偵測器：

\begin{align}
\text{如果它讓我們覺得無聊} &\Rightarrow \text{它會讓讀者睡著},\\
\text{如果它顯得扁平} &\Rightarrow \text{我們就會失去讀者},\\
\text{如果它讓我們喉頭發緊} &\Rightarrow \text{讀者就會落淚},\\
\text{如果它讓我們微笑} &\Rightarrow \text{讀者就會感到溫暖}.
\end{align}

這不是單純的情緒化表述，而是一條可操作規則。Jenkins 的意思是：作者自身切身感受到的反應，本身就是一種證據，說明這段敘事是否仍然帶著力量。

\subsection{Question \& Answer}

\paragraph{Question.} 我們要怎麼知道，哪一個故事點子真的撐得起一部小說？

\paragraph{Answer.} 方法是測它的耐久性。真正該被選中的那個點子，是在中段仍然撐得住我們熱情的那一個；是那個即使工作開始重複、困難、令人洩氣，也仍然把我們拉回鍵盤前的點子。

\paragraph{A worked filter.} 講座中的選擇規則，可以寫成一個短短的更新流程：

\begin{enumerate}
\item 先從一組候選點子開始。
\item 問：哪一個仍然會把我們拉回這項工作？
\item 用小說中段來測它，而不只是用一幅強力的開頭畫面來測它。
\item 保留那個在挫折與重複之下，熱情仍然持續的點子。
\end{enumerate}

\section{氣質、結構與古典序列}

一旦計畫被選出來，講座便追問：我們該如何工作？Jenkins 引入那個大家熟悉的區分：outliner 與 pantser。前者想先把架構鋪好，後者則從一粒種子出發，在寫的過程中發現道路。他還提到 Stephen King，作為第二種類型最知名的例子。

但講座並沒有讓這件事演變成部族之爭。Jenkins 明白地把這種對立消解掉：兩種方法沒有哪一種在本質上更優。許多作家其實都是混合型的，既想要某種地圖帶來的安全，也想要某種即興發現帶來的自由。他自己雖然明顯偏向 pantser 一側，卻立刻補上一句：他從不在完全不知道自己要幹什麼時就開始。這個但書很關鍵，因為它把我們帶向整場講座第一個真正的形式模型。

\subsection{Question \& Answer}

\paragraph{Question.} 我們一定要列大綱嗎？還是可以邊寫邊發現故事？

\paragraph{Answer.} 當然可以在寫的過程中發現許多東西，但發現並不等於結構可以被取消。不論我們的氣質如何，小說仍然需要一條足夠堅固的承重序列，好讓它不會在開頭幾頁之後便停滯。

接著，Jenkins 引入 Dean Koontz 的 ``classic story structure.'' 這是整場講座最清楚的演算法骨架，我們應當把它以緊湊的符號形式保留下來：

\begin{equation}
T_0 \to T_1 \to T_2 \to T_3,
\end{equation}
其中
\begin{align}
T_0&=\text{及早引入的大麻煩},\\
T_1&=\text{不斷惡化的麻煩},\\
T_2&=\text{處境看似已無可挽回},\\
T_3&=\text{主角憑藉發展出的力量與洞見贏得勝利}.
\end{align}

\begin{figure}
\centering
\begin{tikzpicture}
\draw (-1.4,0.55) rectangle (1.4,-0.55);
\node at (0,0.14) {$T_0$};
\node at (0,-0.18) {大麻煩};

\draw (2.6,0.55) rectangle (5.4,-0.55);
\node at (4.0,0.14) {$T_1$};
\node at (4.0,-0.18) {惡化};

\draw (6.6,0.55) rectangle (9.4,-0.55);
\node at (8.0,0.14) {$T_2$};
\node at (8.0,-0.18) {無望};

\draw (10.6,0.55) rectangle (13.4,-0.55);
\node at (12.0,0.14) {$T_3$};
\node at (12.0,-0.18) {掙得的勝利};

\draw[->] (1.4,0) -- (2.6,0);
\draw[->] (5.4,0) -- (6.6,0);
\draw[->] (9.4,0) -- (10.6,0);
\end{tikzpicture}
\end{figure}

這張示意圖不是從影片中擷取出來的視覺證據，而是對 Jenkins 口頭陳述之序列的一次乾淨重建。它所要告訴我們的東西，正是講座所需要的：敘事壓力不是隨機生成的，它是經由早早引入麻煩、逐步加重麻煩、拒絕給出過早舒緩，最後才讓一位已經改變的主角去解決它而被建造出來的。

在結束這一段結構說明之前，Jenkins 給出了一條強烈的命令，也為下一節做好準備：不論我們使用何種基本形式，都必須從一開始就掐住讀者的喉嚨，永遠不要鬆手。直到這之後，他才開始轉向角色。

\section{角色、研究與視角}

講座此時從引擎走向承載者。情節結構本身是空的，除非有角色能夠承受它的壓力。Jenkins 的第一個要求，是主角必須擁有弧線。到了結尾，主要角色必須成為一個不同的人，最好還是更成熟、更好的那個版本。因此，後來的勝利就不可能只是任意的成功，而必須來自某種發展。

第二個條件緊接著出現。主角可以有缺點，但不能厭惡到讓讀者從一開始就被推出去。站在對面，反派則必須強而有力、令人信服。Jenkins 在這裡對初學者的錯誤看得非常清楚：反派不該只因為故事位置上需要一個反派，就成了邪惡的人。大多數反派其實都認為自己是對的。因此，反派同樣需要可理解的動機。

這個原則也向外延伸。那些圍繞主角運轉的重要角色，不能只是來自「中央選角處」的罐頭人物。Jenkins 確實提到還有一張更長的問題清單，用來發展這些人物；但逐字稿在那個關鍵位置剛好破損得很嚴重。因此，我們只保留其中穩定可辨的部分：名字要彼此區分，首字母不要混淆，而角色在外貌與說話方式上，也應該足夠不同，以免讀者把他們搞混。

到了這裡，講座引入一對很有用的耦合變數。主角面對的是一個外部問題，但真正讓他活起來的，是內部動盪。外部目標與內在弱點不能彼此分離；故事的運作，就是把它們逼到一起。

接著，研究進場，而這一部分講座說得非常具體。小說雖然是虛構的，但它仍然必須可信。即使是奇幻與科幻，也必須在內部站得住腳，讀者才會願意暫停不信。Jenkins 給出了一組值得保留的緊湊比例：

\begin{equation}
\text{研究}=\text{調味},\qquad \text{故事}=\text{主菜}.
\end{equation}

然後，他點出幾類一位實際工作的小說家真正可能會用到的資源：

\begin{itemize}
\item 地圖冊與世界年鑑，
\item 百科全書與搜尋引擎，
\item 面對面或通話式訪談，
\item 甚至包括 YouTube 與同義詞詞典這類日常工具。
\end{itemize}

重點不在於光鮮，而在於把地理、文化與技術細節弄對到足以讓讀者不會被甩出故事之外。Jenkins 也補上一個重要但書：在未來題材作品中，我們當然會發明大量技術細節，但即便如此，所有被發明出來的系統仍然必須說得通，並且前後一致。

從研究出發，講座自然轉向 point of view。Jenkins 強調，這不只是語法上第一、第二、第三人稱的選擇而已；它其實是在決定：誰會成為故事的攝影機。令 \(s\) 表示一個場景，\(\mathrm{POV}(s)\) 表示該場景的視角人物，則講座中的操作性限制可以寫成

\begin{equation}
|\mathrm{POV}(s)|=1.
\end{equation}

Jenkins 用散文把這條規則說得非常清楚：一個場景只容許一位視角人物。他接著補充說，自己個人甚至偏好一章一視角，理想上甚至一書一視角。一旦這個上界被設下，敘事資訊的通道也就跟著被限定：

\begin{equation}
\mathrm{Info}(s)=\{\text{看見、聽見、觸到、聞到、嚐到、想到}\}_{\mathrm{POV}(s)}.
\end{equation}

這當然是我們的記號，而不是他的，但它相當忠實地總結了講座的意思：敘述所能傳遞的，是視角角色所感知與所思考的。我們待在那顆心智裡，用那個角色的感官與情緒來過場景。接著，Jenkins 又解開一個初學者常見的誤解。這條限制 \emph{並不} 逼迫我們一定得用第一人稱。第三人稱限知完全可以與單一視角通道相容，而事實上它也是現代敘事中的預設選擇。他真正警告的，是那種隨手而為的全知視角，也就是在同一場景裡不斷跳進跳出不同人物頭腦的老習慣。

\subsection{Question \& Answer}

\paragraph{Question.} 視角究竟如何限制小說可以揭露的東西？

\paragraph{Answer.} 視角是一道資訊邊界，而不是貧乏化。一旦我們選定一個場景的視角人物，散文便會被限制在那個角色所能感知、推知、記起與思考的範圍內。這依然允許我們揭露其他角色，但只能透過一顆意識的壓力來揭露。

\section{從事中開始，並觸發讀者的心智}

就在第六步之前，Jenkins 短暫提醒了聽眾：講座結尾還有那個預告過的 bonus step。這種提醒不是偶然，它讓程序性的框架在講座轉向風格問題時依然保持活著。接著，他以一個古典術語來開啟下一段：\emph{in medias res}，也就是從事中開始。

他很小心地先堵住一個可預期的誤解。從事中開始，並不一定意味著槍戰或追逐。它真正的意思是：當故事本來就應該已經在動時，我們就不從靜止的解說、場景說明或背景鋪墊開始。開頭本身必須帶有向前的壓力。Jenkins 用他一貫乾脆的方式說出規則：每一個句子的目的，甚至每一個字的目的，都是迫使讀者去讀下一句。

我們可以把這條局部運動律寫成

\begin{equation}
\text{第 } n \text{ 句} \Rightarrow \text{想讀第 } n+1 \text{ 句的慾望}.
\end{equation}

接著，他把開場句分成四種類型：驚奇式開場、戲劇性陳述、哲理式開場，以及詩性開場。講座中引用的那些例子，並不是作為文學飾品出場，而是作為不同前進拉力的展示。一口時鐘敲出 \(13\) 下，是對期待的違反；突如其來的暴力，打開了一道戲劇性的缺口；一條哲理性的概括，則邀請我們去驗證它；而一個陌生的詩性圖像，則會使整體語調不穩，逼我們好奇究竟是什麼樣的世界能容納它。

但 Jenkins 並沒有讓講座停留在開場句。強開頭是必要條件，卻不是充分條件。許多小說雖然開頭充滿能量，下一頁卻立刻沉回解說。於是，他接著推進到下一個節點：讀者心智的戲劇場。

這裡的主張是整場講座裡最強的一條之一。作為作者，我們並不是在試圖迫使讀者「剛好以我們看見的方式看見場景」；我們是在試圖觸發他們自己的想像建構。我們只提示到足夠為止，剩下的工作由讀者的心智完成。也就在這裡，講座轉向那條著名的區分：telling 與 showing。

\begin{equation}
\text{告知} \Rightarrow \text{讀者被告知},\qquad
\text{展示} \Rightarrow \text{讀者推知}.
\end{equation}

\subsection{Question \& Answer}

\paragraph{Question.} 所謂「展示，而不是告知」，究竟真正意味著什麼？

\paragraph{Answer.} 它的意思是：我們用可觀察的證據，取代抽象性的報告，好讓讀者從中自行推知那個隱含狀態。讀者不再只是接收一個事實，而是在幫忙建造場景。

\paragraph{A worked update rule.} Jenkins 的轉換程序可以寫成：

\begin{enumerate}
\item 先辨認出那個抽象報告：高、憤怒、寒冷、害怕，等等。
\item 再問：另一個角色，或者讀者，究竟能實際看到什麼？
\item 用這些可觀察線索，取代原本抽象的謂詞。
\item 一旦那些證據已經能自己承載意義，就把那條直接報告拿掉。
\end{enumerate}

講座裡最清楚的兩個例子，可以穩穩地保存為：

\begin{align}
\text{``He was tall''} &\to \text{別人得抬頭看他；他得低頭穿過門框},\\
\text{``He was angry''} &\to \text{他的臉漲紅、喉嚨收緊、聲音揚高、拳頭砸向桌面}.
\end{align}

最後那一步是由讀者在心裡自己完成的。也正因此，展示不只是讓 prose 比較漂亮，而是一種讓讀者參與經驗的方式。

\section{加壓、最黑暗時刻，以及最後的紀律}

到第八步時，前面的結構已經準備好回來，以具體形式再次出現。主角已經被推進麻煩之中；現在的規則則是：每一次試圖逃出去，都應當讓情況更糟。在這裡，講座給出了一句最強的格言：

\begin{equation}
\text{衝突}=\text{小說的引擎}.
\end{equation}

Jenkins 為了強化這個點，故意舉出一個過度舒適的反面例子：一位私家偵探，長得好看、身體健康、有錢、有設備、有幸福家庭、有體面的辦公室，還有一位有錢客戶送來一個 juicy case。除非這是一部諷刺作品，否則這樣的舒適會立刻殺死敘事的力量。他的修正非常粗暴，卻也極其簡單：把支撐拿走。拿掉舒適、金錢、安全、設備、健康、家庭穩定。一旦主角真正受壓，故事才會開始活起來。

我們可以把這條加壓規則重述成一個短推導：

\begin{enumerate}
\item 從真正的麻煩開始。
\item 讓每一個試圖解決的方法，都使情況更糟。
\item 與其增加便利，不如拿掉支撐。
\item 用持續上升的壓力，讓讀者始終待在頁面邊緣。
\end{enumerate}

接著，講座進入第九步，也就是最黑暗的時刻。Jenkins 借用了 Angela Hunt 的說法，用來命名那個連作者自己都會開始懷疑：這故事到底還能怎麼寫下去的時刻。講座中的例子是刻意誇張的：一個曾經浪子回頭、如今看似徹底改過的情人，卻在婚禮前夜親手毀掉了一切。正因如此，這個例子才有用。情境必須看起來真的不可修補。

\begin{definition}
所謂最黑暗的時刻，是情節的局部最低點，是那個連作者都覺得主角的處境已經無望的時刻。
\end{definition}

這於是給出講座後段的那條鏈：

\begin{equation}
\text{麻煩} \to \text{惡化中的麻煩} \to \text{最黑暗時刻} \to \text{高潮} \to \text{結尾}.
\end{equation}

\subsection{Question \& Answer}

\paragraph{Question.} 主角的處境應該絕望到什麼程度，結尾才會顯得是掙來的？

\paragraph{Answer.} 必須絕望到一條簡單出路都看不見。Jenkins 在這裡非常堅決：不要發明一條柔軟的逃生通道，不要臨時塞進一個奇蹟，也不要只因為壓力讓人不舒服，就給主角一個休息。最黑暗時刻的全部功能，就在於逼迫角色只能依賴這個故事一路建構出來的新肌肉與新洞見來行動。

這個回答也為下一個區分鋪平了道路，而 Jenkins 在這裡尤其堅持：

\begin{equation}
\text{高潮} \neq \text{結尾}.
\end{equation}

高潮，是最後的測試，是對抗的峰值，是所有長篇鋪陳終於被支付的點。結尾則發生在那之後。它較安靜，但並非可以省略。Jenkins 對好的結尾提出兩條標準：第一，它把讀者放在第一位，尊重讀者的投入；第二，它會讓英雄一直待在舞台上，直到最後一個字。換言之，小說不能在煙火放完的那一刻就粗暴停止；它還必須真正落地。

講座於是最後一次，從虛構過程回到作者過程。那個 bonus step 並不是附錄，而是在完成整場論證。正如小說需要有序的階段，寫作這件事本身也一樣。Jenkins 最後給出的規則是

\begin{equation}
\text{起草} \neq \text{編輯}.
\end{equation}

第一稿應該在內在批評者被調低音量的情況下寫完。陳腔濫調、粗糙轉場與邏輯上的小問題，不必在當下立刻解決——如果那樣做會讓整體運動停下來的話。修訂應該留到後面，以一種不同的心智模式來完成。因此，講座最後的更新規則就是：

\begin{enumerate}
\item 起草時，不為了完美而停下。
\item 在起草階段允許粗糙存在。
\item 之後，再用獨立的修訂輪回頭處理。
\item 只有當整體的起草運動完成之後，才進行逐字逐詞的微調。
\end{enumerate}

那份預告中的二十一項 checklist，就屬於後面這第二種模式。由於逐字稿沒有保存其具體內容，我們在這裡只保留它的功能：它是修訂工具，而不是起草工具。

\section{Summary}

這場講座從恐懼開始，以紀律收束。一路上，它把一種巨大而模糊的野心，轉化成一條有結構的進程：先界定對象，再選出那個真正能撐過中段的點子，區分氣質與結構，給故事一條承重序列，建出能夠承受壓力的角色，讓世界足夠可信，把散文約束進一條一致的視角，從運動中開始，讓讀者自己推知而不是單純接收，持續增加衝突，把故事一路推到最黑暗的時刻，區分高潮與結尾，最後再把起草與編輯切開。

因此，本章最深的一課並不是說寫作可以被完全約化成公式，而是說：寫作可以被某種公式穩定下來。我們要選擇那個熬得過中段的計畫，要把敘事限制到足以讓讀者真正進入其中，並且讓衝突一路深化，直到那個結尾不只是被抵達，而是真正被掙得。